\chapter{Dental Procedures and Appliances}
\index{Dental Procedures and Appliances}

\section*{Definition and Overview}
Dental appliances are removable or fixed devices used to restore, protect, or straighten the teeth. They include dentures, retainers, orthodontic appliances, mouthguards, bite splints, and space maintainers. Appliances are prescribed to replace missing teeth, correct alignment, protect teeth during sport, manage teeth grinding, or guide dental development in children. They work alongside dental procedures such as impressions, fitting, and adjustments. This chapter explains the common types of dental appliances, what they do, and how to care for them.

\section*{Epidemiology}
Dental appliances are widely used across all ages. Dentures are worn by a large proportion of older people, although rates of complete tooth loss are declining with better prevention. Retainers and orthodontic appliances are used by many teenagers and adults during and after braces. Mouthguards are recommended for all contact sport participants, and bite splints are common for people with teeth grinding, which affects around one in ten adults. Space maintainers help children with early tooth loss preserve space for permanent teeth.

\section*{Aetiology and Pathophysiology}
\begin{itemize}
    \item \textbf{Tooth loss:} dentures, bridges, and implant-supported appliances replace missing teeth
    \item \textbf{Malocclusion:} retainers and orthodontic appliances correct or maintain tooth alignment
    \item \textbf{Injury prevention:} mouthguards absorb impact and protect teeth and jaws during sport
    \item \textbf{Bruxism:} bite splints protect teeth from wear caused by grinding and clenching
    \item \textbf{Growth guidance:} functional appliances and space maintainers guide jaw and dental development in children
    \item \textbf{Mechanism:} appliances apply forces, protect surfaces, or replace missing structures to maintain oral function
\end{itemize}

\section*{Clinical Features}
\begin{itemize}
    \item Dentures restore chewing and appearance after tooth loss
    \item Retainers hold teeth in position after orthodontic treatment
    \item Mouthguards are worn during sport to prevent injury
    \item Bite splints reduce pain and tooth wear from grinding
    \item Space maintainers prevent teeth drifting into a gap after early tooth loss
    \item Appliances need regular cleaning and periodic adjustment or replacement
    \item Discomfort in the first days or weeks after fitting is common and settles
\end{itemize}

\section*{Differential Diagnosis}
\begin{itemize}
    \item Fixed versus removable appliances, depending on the treatment goal
    \item Denture versus bridge versus implant for missing teeth
    \item Retainer versus new orthodontic treatment for drifted teeth
    \item Custom mouthguard versus boil-and-bite versus stock mouthguard
    \item Bite splint versus other treatments for grinding-related pain
\end{itemize}

\section*{When to Seek Medical Care}
See a dentist to discuss appliances for missing teeth, orthodontic needs, sports protection, or grinding. Appliances should be reviewed regularly, and any that cause persistent pain, loosen, or break should be assessed. Children with early tooth loss should have space needs assessed.

\textbf{Contact your local emergency services for:}
\begin{itemize}
    \item A denture or appliance that causes choking or is swallowed
    \item Severe pain, swelling, or bleeding related to an appliance
    \item Injury while wearing a mouthguard that involves the face or head
\end{itemize}

\section*{Diagnosis and Investigations}
\begin{itemize}
    \item \textbf{Dental examination:} determines which appliance is appropriate
    \item \textbf{X-rays:} assess teeth, roots, bone, and space for appliances
    \item \textbf{Impressions or digital scans:} used to fabricate custom appliances
    \item \textbf{Assessment of habits:} grinding, sport participation, and orthodontic needs
    \item \textbf{Fitting and adjustment:} appliances are fitted and adjusted for comfort and function
\end{itemize}

\section*{Management}
\begin{itemize}
    \item \textbf{Dentures:} full or partial, removable replacements for missing teeth, relined or replaced as the mouth changes
    \item \textbf{Retainers:} fixed or removable appliances worn after braces to prevent relapse
    \item \textbf{Orthodontic appliances:} braces, aligners, and functional appliances
    \item \textbf{Mouthguards:} custom-fitted guards provide the best protection and comfort
    \item \textbf{Bite splints:} worn at night to protect teeth from grinding
    \item \textbf{Space maintainers:} fitted in children after early tooth loss
    \item \textbf{Care of appliances:} daily cleaning, safe storage, and regular review
\end{itemize}

\section*{Complications}
\begin{itemize}
    \item Poor fit causing pain, ulcers, or food trapping
    \item Decay or gum disease around fixed appliances if cleaning is inadequate
    \item Breakage, loss, or damage to appliances
    \item Relapse of tooth movement if retainers are not worn
    \item Denture stomatitis (thrush) under ill-fitting dentures
    \item Bone changes under dentures over many years
    \item Costs of replacement and repair
\end{itemize}

\section*{Prognosis and Outcomes}
Dental appliances are highly effective when well made, correctly fitted, and properly cared for. Dentures restore eating and appearance, retainers preserve orthodontic results, mouthguards prevent injury, and splints protect teeth from grinding damage. Regular review and maintenance extend the life of appliances and prevent complications. With good care, most people adapt to their appliances quickly and enjoy lasting improvements in oral health and function.

\section*{Prevention and Self-Care}
\begin{itemize}
    \item Clean removable appliances daily with a soft brush and lukewarm water
    \item Remove appliances as directed and store them safely
    \item Wear mouthguards for all contact sport
    \item Wear splints and retainers exactly as prescribed
    \item Attend regular check-ups so appliances are monitored and adjusted
    \item Replace appliances that no longer fit
    \item Maintain good hygiene around fixed appliances
\end{itemize}

\section*{Sources and Further Reading}
The following reputable sources provide further information for healthcare students and patients seeking reliable, current advice about dental procedures and appliances.

\begin{itemize}
    \item Australian Dental Association. \textit{Dental appliances: information for patients.} https://www.ada.org.au/
    \item NHS. \textit{Dentures, retainers, and mouthguards.} https://www.nhs.uk/conditions/dentures/
    \item Mayo Clinic. \textit{Dentures and other dental appliances.} https://www.mayoclinic.org/
    \item Healthdirect Australia. \textit{Dentures and dental appliances.} https://www.healthdirect.gov.au/
    \item Sports Medicine Australia. \textit{Mouthguards and sports dentistry.} https://sma.org.au/
    \item MedlinePlus. \textit{Dental appliances.} https://medlineplus.gov/
\end{itemize}
